\chapter{The Physics Read Off the Count}

The construction is finished. The machine has been built, climbed, collapsed, needled, walked back into a loop,
counted around, named, and had its split predicted and its operation decomposed. What is left is not more building
but reading: taking the quantities the construction produced and saying, plainly, what physics reads them as. This
chapter is a dictionary. It sets the machine's built and proved quantities in one column and their physical names
in the other, and it does one disciplined thing throughout --- it grades every entry as a reading, a name laid
over a count, never a new theorem. The counts are the construction's; the names are the dictionary's; and the
whole point of having built the counts so carefully is that now, when the physical names are hung on them, you can
see exactly what each name is worth.

\section{The Dictionary}

The claim is that the physics the machine speaks is read off the count --- that each physical quantity the
instrument reports is a name laid over a thing the construction already built or proved.

Take the entries in turn. \emph{Charge} is the loop count: the tally of how many times the circuit was driven
around, from the seam. That is built --- the count is real, kept at the seam --- and ``charge'' is the reading
hung on it. \emph{Mass} is the strain: the second difference that surfaces when the loop is driven past its
threshold, the residue cashed as weight. The surfacing was proved; ``mass'' is the reading of it. \emph{Gravity}
is curvature is the holonomy --- the signed thing the loop accumulates going around, read as the ±1 the last
chapters cashed. Here the dictionary leans on its one genuinely proved leg: that the sign the loop carries flips
two-fold is proved, with no axioms, and ``gravity, the curvature a path accumulates'' is the reading laid over
that proved flip. \emph{Gauge} is the label lines --- the origin and unit and convention the hand hangs on a
reading, the selected part, ours to reset; it is the tange side made physical, and ``gauge'' names it. And
\emph{confinement} --- the fact that certain charges are never seen alone --- is read off the corridor: the band
above the countable where value gives out, where some distinctions are present and cannot be pooled into a count,
shown but not said. What physics calls color and flavor, the things that never appear unbound, the dictionary
reads as the un-poolable residue of that silent band, the difference the machine carries and cannot funge into a
figure.

Now the grade, and it is uniform across the dictionary, which is the whole reason to state it once and clearly.
Every entry above is a \emph{reading} --- an interpretation laid over a built or proved substrate, offered and to
be judged, not itself a theorem. The dictionary does not prove that the loop count is charge, or that the strain
is mass, or that the silent band is confinement; it \emph{reads} them so, and marks each reading as a reading. The
substrate underneath is what carries the weight: the loop count is genuinely built, the strain's surfacing
genuinely proved, the two-fold flip genuinely proved with no axioms, the corridor genuinely constructed as the
band where counting stops. On that built-and-proved substrate the dictionary hangs its physical names, and it
insists you keep the two apart --- the count is the construction's, the physics-word is the reading. There is
exactly one proved leg the dictionary leans on directly, the ±1 of gravity-as-holonomy; every other entry is a
reading over a count, and the honest thing is to say so and not let the dictionary's physical grandeur borrow a
certainty the counts alone possess.

\section{One Fact, Read Many Ways}

Step back from the entries and see what the dictionary is really claiming, because it is bolder as a whole than
any single line. It is claiming that a great swath of physics --- charge, mass, gravity, gauge, the confinement of
the unseen --- is not a heap of separate primitives each demanding its own foundation, but a single thing read
many ways: the count the machine keeps, seen from different sides and given different physical names. Charge and
mass and gravity are not three unrelated quantities that the instrument happened to find; they are the loop count,
the loop's strain, and the loop's holonomy --- three readings of the one closed circuit the construction built.
The physics is unified not by a grand equation but by a common origin: it is all read off the same count, and the
dictionary is the record of the reading.

That is a strong claim and the chapter makes it as a reading, deliberately, because it is exactly the kind of
claim that would be ruinous to overstate. The book does not prove that all of physics reduces to the count; it
reads the physics it has built off the count, entry by entry, and shows that the readings cohere --- that the same
loop yields charge and mass and gravity under three different acts of reading, that the tange side yields gauge
and the silent corridor yields confinement, and that nothing in the dictionary requires a primitive the
construction did not build. What the dictionary demonstrates is coherence, not reduction: that the physics can be
read off the count consistently, each name earning its place over a real construction. Whether that reading is the
\emph{right} one --- whether the world is, at bottom, the count the machine keeps --- is not the dictionary's to
settle and it does not pretend to. It offers the reading, grades it as a reading, and leaves the last word, as
always, to the one holding the book. The physics is read off the count. That the count is the physics, the
dictionary suggests and does not prove.
